\section{Compact spatial realization of the changed two-stage trajectory}
\label{ms:section}

We now realize the actual damped two-stage trajectory, including its
changed first threshold, inherited scale, second threshold, selected
second return, and evolving finite hold, in the full spatial equations.
The coefficients used here are those constructed in the direct return
theorem and Theorem~\ref{ave:complete_entry}. The present calculation supplies their entire
spatial realization; in particular it does not interpret a modal
damping equation as an unforced equation for the released profile.
The physical equations and orientation are
\begin{align*}
 \theta_t+u\cdot\nabla\theta&=d\Delta\theta+f_\theta,\\
 u_t+(u\cdot\nabla)u&=d\Delta u+\theta e_2+f_u,
 &\operatorname{div}u&=0,&p&=0,\\
 J&=\begin{pmatrix}0&-1\\1&0\end{pmatrix},
 &e_2&=(0,1),&x&\in\mathbb R^2.
\end{align*}
Both physical coefficients equal the explicitly selected $d>0$.
The original parameters remain
\[
 A_0\geq1,\qquad\mu=\lambda_0=\lceil\sqrt{A_0}/2\rceil,
 \qquad\nu=\sqrt{A_0},\qquad K_j=\mu\lambda_j\quad(j=1,2).
\]
Thus $d$ is not the source time-scaling parameter $\nu$.
Here $\lambda_j$ denotes exactly the source frequency written
$\widehat\lambda_j$ in the preceding entry calculation.

\subsection{The actual coefficients and the nested material regions}

Write $e(s)=(\sin s,\cos s)$, $e_0=R_\alpha e_2$,
$p_1=e(s_*)$, and $p_2=e(s)$, where
$s=L_2\nu/\sigma$, $a=\sin s$, $h_0=\cos s$, and
$\Gamma=\sigma a$ use the actual diffusive inherited scale $\sigma$.
The activation functions are
\[
 w_1=\eta(\Gamma_1(t-t_{01})),\qquad
 w_2=\eta(\Gamma(t-t_{02})),\qquad
 \Gamma_1=\nu\sin s_*,\qquad t_{01}=1/\nu.
\]
The source step $\eta$ is smooth nondecreasing, zero on
$(-\infty,0]$ and one on $[1,\infty)$.
All inactive increments below are defined to be zero. The finite
coefficient construction has the exact equations
\begin{align}
 G_{<j}&=-A_0e_0+\sum_{i<j}w_iK_i\Theta_i\zeta_i,
 &D_{<j}&=\dot\alpha J+
       \sum_{i<j}\frac{w_i\Omega_i}{r_i^2}J\zeta_i\otimes\zeta_i,
 \label{ms:parents}\\
 \dot M_j&=D_{<j}M_j,&M_j(t_{0j})&=I,
 &\zeta_j&=M_j^{-T}p_j,&r_j^2&=|\zeta_j|^2,\nonumber\\
 \dot\Theta_j&=-\frac{J\zeta_j\cdot G_{<j}}{K_jr_j^2}\Omega_j
                   -dK_j^2r_j^2\Theta_j,
 &\dot\Omega_j&=K_j\zeta_{j,1}\Theta_j-dK_j^2r_j^2\Omega_j.
 \label{ms:amplitudes}
\end{align}
The common angle is the selected original first control followed by
the exact second feedback control. In particular the first threshold
time is $t_{01}+L_1/(\Gamma_1-dK_1^2)$, not the inviscid
threshold time. The source first transition and pulse are translated
to that actual threshold. During second growth and transition,
$\alpha=s_*$, $M_2=I$, $\Omega_1=0$, and $\Theta_1$ diffuses.
The second activation time is exactly $t_{02}=t_{b1}$, the actual
first return time, with no intervening gap.
During second steering and holding the exact inverse is
\begin{align}
 M_1&=R_\alpha,&
 M_2&=R_{\alpha-s_*}
       \begin{pmatrix}1&-(h-h_0)/a\\0&1\end{pmatrix},
 &\Omega_1&=\sigma W,&\Theta_1&=-\sigma^2B/K_1,
 \label{ms:M}\\
 h'&=W,&
 W'&=B\frac{Y-hz}{h^2+a^2}-\delta_1W,&
 B'&=-C_1W-\delta_1B,&
 Y&=\sqrt{h^2+a^2-a^2z^2},\nonumber
\end{align}
where $\tau=\Gamma(t-t_1)$, a prime denotes $d/d\tau$,
$\delta_1=dK_1^2/\Gamma$,
$C_1=A_0\sin s_*/(\sigma^2a)$, and $z$ is the selected
source ramp, pulse, and zero hold. The proved bounds give
\[
 h_0\leq h<3,\quad 0\leq W<9,\quad
 1/4<B\leq2,\quad h_0\geq\sqrt{63}/8,\quad
 r_1=1,\quad h_0\leq r_2<\sqrt{10}.
\]
These are actual changing coefficients, not independently chosen
affine backgrounds. Every $D_{<j}$ has trace zero: the rotation
has trace zero and
$\operatorname{tr}(J\zeta_i\otimes\zeta_i)=\zeta_i\cdot J\zeta_i=0$.
Jacobi's determinant formula therefore gives $\det M_j=1$.
The phase equation follows by differentiating the inverse transpose:
$\dot\zeta_j=-D_{<j}^T\zeta_j$.

On $[-\pi,\pi]$ define the released profile by the following formula
and then extend it periodically to $\mathbb R$:
\[
 F(s)=\sin s+\chi(s)(s-\sin s),\qquad
 P'=F,\qquad\int_{-\pi}^{\pi}P(s)\,ds=0.
\]
Here $\chi$ is even, supported in $(-\pi,\pi)$, and equals one
on $[-s_0,s_0]$, with $0<s_0<1$. Oddness and periodicity give
a unique even periodic primitive $P$ with the specified mean.
In particular $F(s)=s$ on $|s|\leq s_0$; its primitive there is
$P(0)+s^2/2$, with its original constant $P(0)$ retained.
Let $f\in C_c^\infty(B_1)$ be radial, $0\leq f\leq1$,
and $f=1$ on $B_{1/2}$. Retain the source radii
\[
 \ell_1=\frac{s_0}{C_\ell\mu},\qquad
 \ell_2=\frac{\lambda_1^{-3}}\mu,\qquad C_\ell\geq4.
\]
Define the material coordinates and envelopes
\[
 y_j=M_j^{-1}x,\qquad g_j(x,t)=f(y_j/\ell_j),\qquad
 s_j=K_j\zeta_j\cdot x.
\]
With $N=1+(3-h_0)/a$, the matrix expression \eqref{ms:M}
and $\|I+qE_{12}\|\leq1+|q|$ show
\begin{equation}\label{ms:N}
 \|M_1\|=\|M_1^{-1}\|=1,
 \qquad\|M_2\|,\|M_2^{-1}\|\leq N.
\end{equation}
The second inequality also holds before steering, where $M_2=I$.
The entry construction, \eqref{ave:nesting}, selects $\lambda_1$ with the explicit strict
inequalities
\begin{equation}\label{ms:nesting}
 N\lambda_1^{-3}<\frac{s_0}{2C_\ell},
 \qquad N\lambda_1^{-2}<s_0.
\end{equation}
Consequently every point of the second support satisfies
$|x|\leq N\ell_2<\ell_1/2$ and
$|K_1\zeta_1\cdot x|\leq K_1N\ell_2<s_0$.
It therefore lies in the exact affine cell of the first wave for
the whole construction. The first support lies in $B_{\ell_1}$;
since $\mu\ell_1=s_0/C_\ell<s_0$, the base is exactly affine
there for every angle. Strictness leaves an open neighborhood of
each support on which the corresponding affine identities hold.

\subsection{Original datum, fields, and complete residual forces}

Retain the original radial cutoff $\chi_0\in C_c^\infty(B_{R_0})$,
equal to one on $B_2$, and the exact datum
\[
 \theta_{\rm in}(x)=-\frac{A_0}{\mu}\sin(\mu x_2)\chi_0(x),
 \qquad u_{\rm in}=0.
\]
Choose the radial $H_0\in C_c^\infty(B_{R_H})$ used in the
base construction, where $R_H>R_0+1$ and
$H_0=|x|^2/2$ on $B_{R_0+1}$. Put
\begin{align}
 \theta_0&=-\frac{A_0}{\mu}F(\mu e_0\cdot x)\chi_0(x),
 &\theta_{\rm ref}&=-\frac{A_0}{\mu}F(\mu x_2)\chi_0(x),\nonumber\\
 \theta_b&=\theta_0+(1-\eta(\nu t))
                          (\theta_{\rm in}-\theta_{\rm ref}),
 &\psi_b&=\dot\alpha H_0,&u_b&=J\nabla\psi_b.
 \label{ms:base}
\end{align}
The angle is zero before $1/\nu$; the initial interpolation is
therefore exact and the initial velocity is zero. Define
\begin{equation}\label{ms:fields}
 \begin{gathered}
 T_j=w_j\Theta_j,\qquad Q_j=\frac{w_j\Omega_j}{K_j^2r_j^2},
 \qquad V_j=T_jF(s_j)g_j,\qquad\Psi_j=Q_jP(s_j)g_j,\\
 v_j=J\nabla\Psi_j,\qquad
 \theta=\theta_b+V_1+V_2,\qquad u=u_b+v_1+v_2.
 \end{gathered}
\end{equation}
The flat activation functions make the zero extensions of the
increments smooth, with all spatial and temporal derivatives, at
$t=t_{0j}$. The velocity is divergence free because mixed partial
derivatives commute.

Hereafter fix a layer and temporarily suppress its index. Write
$k=K\zeta$, $\rho=|k|^2=K^2r^2$, $D=D_{<j}$, $G=G_{<j}$,
and $b=k_1$. All functions of one variable in the formulas below
are evaluated at $s=k\cdot x$. Define
\begin{equation}\label{ms:Qrest}
 Q_* =\frac{\dot w\Omega+wb\Theta}{\rho}
                         -Q\frac{\dot\rho}{\rho},
 \qquad\dot\rho=-2k\cdot D^Tk.
\end{equation}
Thus $\dot Q=Q_*-d\rho Q$, directly from
\eqref{ms:amplitudes}, without a constant-phase-length assumption.
The complete added scalar and vector forces are
\begin{align}
 E_\theta={}&\dot w\Theta Fg+QP(J\nabla g\cdot G)
       +QT(F^2-PF')g(Jk\cdot\nabla g)\nonumber\\
 &\quad-dT\{\rho(F+F'')g+2F'k\cdot\nabla g+F\Delta g\},
 \label{ms:scalar}\\
 E_u={}&2Dv+Q_*J\{kFg+P\nabla g\}
                         +(\nabla v)v-TFg e_2
                         -d(\rho v+\Delta v).
 \label{ms:vector}
\end{align}
These are full vector forces at the declared pressure $p=0$.
For completely componentwise evaluation one uses
\begin{align}
 v&=QJ\{kFg+P\nabla g\},\nonumber\\
 \nabla v&=QJ\{k\otimes k F'g+
           F(k\otimes\nabla g+\nabla g\otimes k)+P\nabla^2g\},
 \label{ms:velocityjets}\\
 \Delta V&=T\{\rho F''g+2F'k\cdot\nabla g+F\Delta g\},
 \label{ms:lapscalar}\\
 \Delta v&=QJ\{\rho kF''g+\rho F'\nabla g
       +2F'k(k\cdot\nabla g)+2F(\nabla^2g)k
                         +Fk\Delta g+P\nabla\Delta g\}.
 \label{ms:lapvector}
\end{align}
Thus every phase, envelope, and mixed Laplacian term is displayed.
The base forces are
\begin{align}
 f_{\theta,b}&=\nu\eta'(\nu t)
                 (\theta_{\rm ref}-\theta_{\rm in})-d\Delta\theta_b,
 \nonumber\\
 f_{u,b}&=\ddot\alpha J\nabla H_0+
        \dot\alpha^2\bigl(\nabla(J\nabla H_0)\bigr)J\nabla H_0
                          -\theta_b e_2-d\dot\alpha J\nabla\Delta H_0.
 \label{ms:baseforce}
\end{align}
The requested total forces are exactly
\begin{equation}\label{ms:totalforce}
 f_\theta=f_{\theta,b}+E_{\theta,1}+E_{\theta,2},
 \qquad f_u=f_{u,b}+E_{u,1}+E_{u,2}.
\end{equation}

\begin{proof}[Proof of the full spatial equations]
For the material derivative $\mathcal D=\partial_t+Dx\cdot\nabla$,
differentiating $M^{-1}x$ gives $\mathcal D(M^{-1}x)=0$.
The phase equation gives $\mathcal D(k\cdot x)=0$ as well.
Consequently $\mathcal Dg=0$, $\mathcal DV=\dot T Fg$,
and $\mathcal D\Psi=\dot QPg$.
On this support the actual older temperature and velocity are
$G\cdot x$ and $Dx$ by \eqref{ms:nesting}; in particular their
gradient and Jacobian are exactly $G,D$, even while these change.
The temperature advection increment is therefore
$\dot T Fg+v\cdot G+v\cdot\nabla V$.
Using $Jk\cdot k=0$ and $J\nabla g\cdot\nabla g=0$ gives
\[
 v\cdot\nabla V=QT(F^2-PF')g(Jk\cdot\nabla g),\qquad
 v\cdot G=QFg(Jk\cdot G)+QP(J\nabla g\cdot G).
\]
Since $\dot\Theta=-(Jk\cdot G)\Omega/\rho-d\rho\Theta$,
the first term in $v\cdot G$ cancels the coupled part of
$\dot T Fg$ exactly. Subtraction of $d\Delta V$, computed
by the product rule, proves \eqref{ms:scalar}.

For every trace-free two by two matrix, direct multiplication
gives $DJ+JD^T=0$. Differentiation of a gradient under $\mathcal D$
therefore gives
\[
 \mathcal D(J\nabla\Psi)=D J\nabla\Psi+
                                  J\nabla(\mathcal D\Psi).
\]
The other linear advection increment is
$(v\cdot\nabla)(Dx)=Dv$. Adding it to this formula, adding the
self-advection $(\nabla v)v$, and subtracting $Ve_2+d\Delta v$
proves \eqref{ms:vector}, after using $\dot Q=Q_*-d\rho Q$.
Differentiating $\Delta\Psi=Q(\rho F'g+2Fk\cdot\nabla g+P\Delta g)$
and applying $J$ proves all six groups in \eqref{ms:lapvector}.

The base phase $\mu e_0\cdot x$ is transported by
$\dot\alpha Jx$, and its radial cutoff is also transported.
Before $1/\nu$ its velocity is zero. This proves the scalar
base formula. Differentiating $\dot\alpha J\nabla H_0$ proves
the vector base formula directly on all of $\mathbb R^2$.
Finally, adding increments in the order $1,2$ accounts exactly
for both directions of each interaction. In the second increment,
$\mathcal D$ contains $u_b+v_1$, and the terms $v_2\cdot G_{<2}$
and $D_{<2}v_2$ contain the reverse interactions with that full
older state. Outside an increment's support all its derivatives
vanish. Thus the residual identities telescope globally to
\eqref{ms:totalforce}, with no omitted interaction and no inverse
curl or pressure reconstruction.
\end{proof}

The material-coordinate form makes the original metric equally
explicit. Set $C=M^{-1}M^{-T}$ and $\widetilde g(y)=f(y/\ell)$.
The pullback of the spatial derivative and Laplacian on a function
of $(s,y)$ are
\begin{equation}\label{ms:operators}
 \mathcal P=K\zeta\partial_s+M^{-T}\nabla_y,
 \qquad\mathcal L=K^2r^2\partial_s^2
             +2K(Cp\cdot\nabla_y)\partial_s+
                       \sum_{l,m}C_{lm}\partial_{y_l}\partial_{y_m}.
\end{equation}
Indeed $s=Kp\cdot y$, $\nabla_x=M^{-T}\nabla_y$ on a
material function, and $|M^{-T}p|^2=p^TCp=r^2$; multiplying
the two derivative operators proves \eqref{ms:operators}.
For example $k\cdot\nabla_xg=K(Cp)\cdot\nabla_y\widetilde g$.
This proves an exact coordinate correspondence for every displayed
mixed derivative and keeps the source radii in their original units.

\subsection{What the affine cells and sine regions actually satisfy}

On $g=1$ and $|s|<s_0$ the wave is exactly
\[
 V=T k\cdot x,\qquad
 v=QJk\otimes k\,x
   =\frac{w\Omega}{r^2}J\zeta\otimes\zeta\,x.
\]
Both Laplacians vanish, but their damping contributions to the
forces \eqref{ms:scalar}--\eqref{ms:vector} are
\begin{equation}\label{ms:coreforces}
 -d\rho V,\qquad -d\rho v.
\end{equation}
They do not vanish unless the corresponding wave vanishes.
For the first wave the scalar core force is
$-dK_1^2(w_1K_1\Theta_1\zeta_1)\cdot x$ after activation,
and the curl of its damping vector force is
$-dK_1^2w_1\Omega_1$. Thus the evolving affine parent itself
has the force required by its changed amplitude equations.
On an uncut region where $F(s)=\sin s$, instead,
$F''=-F$ and the two damping contributions are zero:
$\rho V+\Delta V=0$ and $\rho v+\Delta v=0$.
The non-damping vector terms in \eqref{ms:vector} are still
present; this cancellation is not a claim of unforced momentum
at $p=0$. In cutoff regions every term in
\eqref{ms:lapscalar}--\eqref{ms:lapvector} must also be kept.
These restrictions prove exactly how the damped coefficient
dynamics map into the released spatial profile. They differ
from retaining an inviscid state and adding only $-d\Delta$ of
that state: the present amplitudes, parent, phases and controls
have all changed according to \eqref{ms:amplitudes}.

\subsection{Explicit bounds at every spatial derivative order}

For a scalar, vector, or tensor field $U$, let
$|U|_m=\sup_x\|D_x^mU(x)\|$, using Euclidean multilinear
operator norms and $m\geq0$. For a one-variable profile $W$ write
$[W]_m=\sup_s|W^{(m)}(s)|$, and put $F_m^f=\sup_y\|D^mf(y)\|$.
Set $N_1=1$, $N_2=N$, $r_{1,-}=r_{1,+}=1$,
$r_{2,-}=h_0$, $r_{2,+}=\sqrt{10}$. Define finite explicit
numbers
\begin{align}
 \kappa_j^*&=K_jr_{j,+},&
 e_{j,m}&=(N_j/\ell_j)^m F_m^f,\nonumber\\
 \mathcal B_{j,m}(W)&=\sum_{q=0}^m\binom mq
       (\kappa_j^*)^q[W]_q e_{j,m-q},&
 \mathcal C_{j,m}(P)&=\sum_{q=0}^m\binom mq
       (\kappa_j^*)^q[P]_q e_{j,m-q+1}.
 \label{ms:Bnorm}
\end{align}
Here $\kappa_j^*$ is merely a frequency bound, not a diffusivity.
Repeated differentiation and the chain rule give
$|g_j|_m\leq e_{j,m}$,
$|W(s_j)g_j|_m\leq\mathcal B_{j,m}(W)$, and
$|P(s_j)\nabla g_j|_m\leq\mathcal C_{j,m}(P)$.
To prove the middle inequality, assign each of the $m$ derivatives
either to the profile or to the envelope. There are $\binom mq$
assignments with $q$ profile derivatives, each bounded by the
corresponding summand. The other inequalities follow by the same
argument. This proof applies to all directions in the defined
operator norm, not only coordinate derivatives.

The actual threshold amplitudes are
$P_j^*=(\nu^2/\mu)\lambda_j^{-7/8}$. The finite coefficient
calculation, \eqref{ave:amplitude_sup}, supplies the following global
bounds through the hold:
\begin{align}
 U_{\Theta,1}&=\max\{e^{2+\Lambda_1^{-1}}P_1^*,2\sigma^2/K_1\},
 &U_{\Omega,1}&=\max\{(K_1/\nu)e^{2+\Lambda_1^{-1}}P_1^*,9\sigma\},
 \nonumber\\
 U_{\Theta,2}&=e^{18}P_2^*,
 &U_{\Omega,2}&=3K_2U_{\Theta,2}/\sigma.
 \label{ms:amplitudebounds}
\end{align}
In particular these preserve the actual source frequencies and
the changed inherited scale. Put
\[
 A_1=\sup_{[0,T_f]}|\dot\alpha|,\quad
 A_2=\sup_{[0,T_f]}|\ddot\alpha|,\quad
 D_1^*=A_1,\quad D_2^*=A_1+U_{\Omega,1},\quad
 G_1^*=A_0,\quad G_2^*=A_0+K_1U_{\Theta,1}.
\]
Explicit profile bounds for $A_1,A_2$ are given below. Set
$W_{1,1}=\Gamma_1[\eta']_0$ and
$W_{2,1}=\Gamma[\eta']_0$. Then
\begin{align}
 Q_j^*&=\frac{U_{\Omega,j}}{K_j^2r_{j,-}^2},\nonumber\\
 Q_{*,j}^*&=\frac{W_{j,1}U_{\Omega,j}
          +\kappa_j^* U_{\Theta,j}+2D_j^*U_{\Omega,j}}
                       {K_j^2r_{j,-}^2}.
 \label{ms:Qbounds}
\end{align}
The bounds $|Q_j|\leq Q_j^*$ and $|Q_{*,j}|\leq Q_{*,j}^*$
follow from \eqref{ms:Qrest} and
$|\dot\rho/\rho|\leq2\|D\|$. Define
\begin{align}
 \mathcal D_{\theta,j,m}
 &=dU_{\Theta,j}\{(\kappa_j^*)^2\mathcal B_{j,m}(F)
                                      +2\mathcal B_{j,m+2}(F)\},\nonumber\\
 \mathcal D_{u,j,m}
 &=dQ_j^*\{(\kappa_j^*)^2\mathcal B_{j,m+1}(P)
                                      +2\mathcal B_{j,m+3}(P)\}.
 \label{ms:diffnorm}
\end{align}
These bound the complete damping-plus-diffusion terms in
\eqref{ms:scalar}--\eqref{ms:vector}. In particular they include
the core damping and every envelope derivative. The two in front
of $\mathcal B_{m+2}$ or $\mathcal B_{m+3}$ is the sum of the
two coordinate Laplacian terms; $J$ has operator norm one.
The whole added forces satisfy the explicit inequalities
\begin{align}
 |E_{\theta,j}|_m\leq{}&W_{j,1}U_{\Theta,j}\mathcal B_{j,m}(F)
   +Q_j^*G_j^*\mathcal C_{j,m}(P)+\mathcal D_{\theta,j,m}\nonumber\\
 &+Q_j^*U_{\Theta,j}\sum_{q=0}^m\binom mq
        \mathcal B_{j,q+1}(P)\mathcal B_{j,m-q+1}(F),
 \label{ms:fullscalarnorm}\\
 |E_{u,j}|_m\leq{}&(2D_j^*Q_j^*+Q_{*,j}^*)\mathcal B_{j,m+1}(P)
                  +U_{\Theta,j}\mathcal B_{j,m}(F)+\mathcal D_{u,j,m}
 \nonumber\\
 &+(Q_j^*)^2\sum_{q=0}^m\binom mq
            \mathcal B_{j,q+1}(P)\mathcal B_{j,m-q+2}(P).
 \label{ms:fullvectornorm}
\end{align}
For the scalar nonlinear term it is enough to use its exact
alternative expression $v_j\cdot\nabla V_j$; differentiation
of this product proves the displayed sum. For the vector nonlinear
term differentiate $(\nabla v_j)v_j$ and interchange the summation
index to obtain the final sum. These arguments prove all four
bounds directly from the exact formulas.

The base has equally explicit bounds. Put
$c_m=|\chi_0|_m$, $h_m=|H_0|_m$, and
\[
 b_m(W)=\sum_{q=0}^m\binom mq\mu^q[W]_q c_{m-q},
 \qquad B_m^0=\frac{A_0}{\mu}\max\{b_m(F),b_m(\sin)\}.
\]
The interpolation is a convex combination when $\alpha=0$;
afterwards $|\mu e_0|=\mu$. Hence $|\theta_b|_m\leq B_m^0$.
The product rule and \eqref{ms:baseforce} give
\begin{align}
 |f_{\theta,b}|_m\leq{}&
       \frac{\nu A_0}{\mu}[\eta']_0\{b_m(F)+b_m(\sin)\}+2dB_{m+2}^0,
 \nonumber\\
 |f_{u,b}|_m\leq{}&A_2h_{m+1}+B_m^0+2dA_1h_{m+3}
       +A_1^2\sum_{q=0}^m\binom mq h_{q+2}h_{m-q+1}.
 \label{ms:basenorm}
\end{align}
All total force norms are bounded by adding their base and two
increment bounds. Integrating any bound over its interval multiplies
its constant by that interval's physical length. For example
\begin{equation}\label{ms:timecost}
 \int_0^{T_f}|\hbox{wave damping-plus-diffusion force}_j|_m\,dt
 \leq (T_f-t_{0j})\mathcal D_{\bullet,j,m}.
\end{equation}
Here $T_f=t_1+(1+L_2^{-1}+1/32)/\Gamma$ is the actual end
of the finite hold and the dot denotes the corresponding scalar or
vector bound. Every frequency, envelope radius and diffusion
coefficient is explicit. No estimate in this section asserts a bound
uniform over a sequence $K_j\to\infty$.

\subsection{Control derivatives and every mixed derivative order}

The constants $A_1,A_2$ in the spatial force costs can be bounded
without differentiating an unknown shooting root. The root is a
fixed number in $[0,3]$, and differentiation below is in time.
For a compression $L$ put
\[
 Z_0(L)=\max\{1,3L[\eta_2]_0\},\qquad
 Z_b(L)=\max\{[\eta^{(b)}]_0,
                         3L^{b+1}[\eta_2^{(b)}]_0\}\quad(b\geq1).
\]
These bound the derivatives of the source ramp and pulse $z$;
the hold has zero derivatives. During the first control,
$\alpha=s_*-\arcsin(\sin s_*z)$ and $|\sin s_*z|\leq1/4$.
Twice differentiating this displayed formula gives
\begin{align*}
 A_{1,\mathrm{first}}&=\frac4{\sqrt{15}}\Gamma_1\sin s_*Z_1(\Lambda_1),\\
 A_{2,\mathrm{first}}&=\Gamma_1^2\left\{
   \frac4{\sqrt{15}}\sin s_*Z_2(\Lambda_1)
  +\frac{64}{15\sqrt{15}}\sin^3s_*Z_0(\Lambda_1)Z_1(\Lambda_1)^2
                                                    \right\}.
\end{align*}
During second steering and holding write $z_b=Z_b(L_2)$ and
$y_-=\sqrt{59}/8$. The exact control is
\[
 \alpha_\tau=\frac{W\zeta_{1,1}-az'}Y,
 \qquad Y'=\frac{hW-a^2zz'}Y,
 \qquad\zeta_{1,\tau}=\alpha_\tau J\zeta_1.
\]
Define the explicit constants
\begin{align*}
 U_\alpha&=(9+az_1)/y_-,\qquad
 U_Y=(27+az_1/4)/y_-,\\
 U_{W'}&=\frac{128}{63}(\sqrt{10}+3z_0)+\frac9{16},\\
 U_{\alpha'}&=(U_{W'}+9U_\alpha+az_2)/y_-
                                      +(9+az_1)U_Y/y_-^2.
\end{align*}
The bounds $Y\geq y_-$, $|az|\leq1/4$, $h<3$, $W<9$,
$B\leq2$, $h^2+a^2\geq63/64$, and $\delta_1\leq1/16$
give $|\alpha_\tau|\leq U_\alpha$, $|Y'|\leq U_Y$ and
$|W'|\leq U_{W'}$. Differentiating the quotient for
$\alpha_\tau$ gives $|\alpha_{\tau\tau}|\leq U_{\alpha'}$.
Thus one can use
\[
 A_1\leq\max\{A_{1,\mathrm{first}},\Gamma U_\alpha\},
 \qquad A_2\leq\max\{A_{2,\mathrm{first}},\Gamma^2U_{\alpha'}\}.
\]
On intervening growth and transition intervals the angle is constant.

For completeness the following exact finite recursion evaluates any
mixed derivative, rather than treating spatial estimates as time
estimates. For a fixed layer put $A=M^{-1}/\ell$ and introduce
independent variables $t,x,s,y$. Define operators on smooth
expressions in these variables by
\begin{equation}\label{ms:jetops}
 \mathfrak T=\partial_t+(\dot k\cdot x)\partial_s
                          +(\dot A x)\cdot\nabla_y,
 \qquad
 \mathfrak X_i=\partial_{x_i}+k_i\partial_s+
                              \sum_l A_{li}\partial_{y_l}.
\end{equation}
Here $\partial_t$ differentiates every time-dependent coefficient
and $\partial_{x_i}$ differentiates the explicitly occurring $x$.
If $\iota^*R=R(t,x,k(t)\cdot x,A(t)x)$, the chain rule proves
\begin{equation}\label{ms:jetpullback}
 \partial_t^b\partial_x^\gamma(\iota^*R)
       =\iota^*(\mathfrak T^b\mathfrak X_1^{\gamma_1}
                                         \mathfrak X_2^{\gamma_2}R).
\end{equation}
Induction proves this identity for all nonnegative $b,\gamma_i$.
For example, for $R=Q(t)P(s)f(y)$ it gives every streamfunction
jet; applying one more $\mathfrak X$ and $J$ gives every velocity
jet. The commutators vanish: in $[\mathfrak T,\mathfrak X_i]$
the derivative $\dot k_i$ is canceled by differentiating
$\dot k\cdot x$, and $\dot A_{li}$ by differentiating
$(\dot A x)_l$. The $\mathfrak X_i$ commute by their displayed
coefficients. This checks the order consistency of the recursion.

Apply \eqref{ms:jetpullback} to each summand of
\eqref{ms:scalar}--\eqref{ms:baseforce}, using a separate pair
$(s,y)$ for each layer when necessary. After expansion the answer
is a finite sum of terms of the form
\[
 c(t)x^\beta\prod_{l=1}^q W_l^{(n_l)}(s_j)
                         \prod_{l=1}^r\partial^{\gamma_l}f(y_j),
\]
or the analogous base terms with $\chi_0,H_0$. An explicit bound
for each term is obtained by replacing it by
\[
 \sup_{[0,T_f]}|c(t)|\,(N_j\ell_j)^{|\beta|}
          \prod_{l=1}^q[W_l]_{n_l}
          \prod_{l=1}^r F^f_{|\gamma_l|},
\]
and summing; for a base term use $R_H^{|\beta|}$.
Each coefficient jet in $c$ is computed by differentiating
\eqref{ms:amplitudes}, the phase equations, \eqref{ms:M}, and
the explicit controls. This is a finite recursion using only
addition, multiplication, the stated profile derivatives and the
nonzero denominators $a$, $K_j$, $\Gamma$, $r_j^2$, $Y$.
Their positive lower bounds have already been proved. Induction
on the requested order, applying the product and quotient rules
to those equations, therefore supplies finite bounds for every
coefficient jet from the preceding-order bounds. This specifies
every mixed force and diffusion norm by an exact finite procedure,
including the time variation of the material cutoff. It does not
replace the cutoff by a stationary radial function for layer two.

\subsection{Smooth finite slab, compact forces, and exact returned state}

All fields and forces are smooth and supported in $B_{R_H}$ on
$0\leq t\leq T_f$. Each velocity is the rotated gradient of a
compact streamfunction, so its energy is finite. It also has the
usual exact Biot--Savart relation. Indeed for
$N_0(x)=(2\pi)^{-1}\log|x|$ and compact smooth
$\psi=\psi_b+\Psi_1+\Psi_2$, distributional integration by
parts gives $N_0*\Delta\psi=(\Delta N_0)*\psi=\psi$;
applying $J\nabla$ gives $u=J\nabla N_0*\omega$ with
$\omega=\Delta\psi$. All integrations are legitimate because
$\psi$ is compact and $N_0$ is locally integrable.
Odd $F$, even $P$, and even cutoffs give odd $\theta,u,f_\theta,f_u$
and even $\omega$. The pressure $p=0$ has $p(0,t)=0$.

At the selected second return and throughout its physical hold
of duration $1/(32\Gamma)$, the coefficient calculation proves
$\zeta_{2,1}=\Omega_2=0$, $w_2=1$, and
\[
 \Theta_2(t)=-P_2(t_b)
     \exp\left[-dK_2^2\int_{t_b}^t|\zeta_2(r)|^2\,dr\right].
\]
Thus $v_2=0$ everywhere on this hold, while its temperature and
the first parent continue changing. In the inner second affine
cell the total fields have exactly
\[
 \theta=\left(-A_0e_0+K_1\Theta_1\zeta_1
                          +K_2\Theta_2\zeta_2\right)\cdot x,
 \qquad
 u=\left(\dot\alpha J+\Omega_1J\zeta_1\otimes\zeta_1\right)x.
\]
The second holding shear is zero, but the first holding shear is
positive at return and evolves by the displayed system. In particular
these exact full spatial fields have the origin vorticity calculated
in \eqref{sr:origin_vorticity}:
\begin{equation}\label{ms:originvorticity}
 \operatorname{curl}u(0,t)=2\dot\alpha+\Omega_1
 =\Omega_1\left(1-\frac{2a^2}{h^2+a^2}\right)
 \geq\frac{31}{32}\Omega_1>0\qquad(t_b\leq t\leq T_f).
\end{equation}
Indeed on holding $\alpha=s_*+\arctan(a/h)$ and
$\dot h=a\Omega_1$, so $\dot\alpha=-a^2\Omega_1/(h^2+a^2)$.
Moreover $h\geq h_0$, $h_0^2+a^2=1$ and $a\leq1/8$ prove
the numerical lower bound. Positivity of $\Omega_1$ at return
persists because the holding equation has positive forcing $B/r_2$.
These are the actual final states realized by the full spatial equations.

There is also an honest compact extension of the forces in time.
For $t<0$ extend the initial state by its constant values, with
$\alpha=0$ and both increments zero. Flatness of $\eta$ makes
the formulas smooth at zero. To the right, continue $z=0$ and
the exact holding ODE for a collar of physical length
$\epsilon_R=1/(128\Gamma)$. Its extension is quantitative:
the total scaled duration is at most $17/16$ since $L_2\geq64$;
the same integrating-factor estimates give
\[
 0\leq W\leq26/3,\qquad
 h\leq1+16/15+(26/3)(7/128)=813/320<3,
 \qquad B>191/512>1/4.
\]
For the last bound use $B_{\rm in}\geq1/2$,
$\delta_1\leq1/16$, $C_1\leq1/32$ and
$B=e^{-\delta_1\tau}B_{\rm in}
       -C_1\int_0^\tau e^{-\delta_1(\tau-r)}W(r)dr$.
These strict bounds close the positivity bootstrap and keep all
denominators away from zero; the smooth finite-dimensional vector
field therefore continues through this collar. The support bounds
continue as well. Extend the newest temperature by its exact hold
exponential, and the angle and matrices by their exact inverse.
In particular no nonstationary endpoint is frozen.

Choose $\epsilon_L=1/\nu$ and multiply the extended residual
forces, not the fields on their required slab, by
\[
 \beta(t)=\eta(1+t/\epsilon_L)
                       \eta(1+(T_f-t)/\epsilon_R).
\]
This function is one on $[0,T_f]$, zero outside
$(-\epsilon_L,T_f+\epsilon_R)$, and flat at the two external
endpoints. The resulting forces belong to
$C_c^\infty(\mathbb R^2\times\mathbb R)$ and still give exactly
the solution already proved on $[0,T_f]$. On the collar the same
mixed-derivative recursion applies, and the product rule includes
the explicit additional factors $\epsilon_L^{-b}[\eta^{(b)}]_0$
and $\epsilon_R^{-b}[\eta^{(b)}]_0$. No equation for the
continued fields outside the required slab is claimed after this
force cutoff. This completes the compact finite spatial realization
and its full derivative and support costs.
